% ═══════════════════════════════════════════════════════════════
% LERNBLATT: Tabellenkalkulation Grundlagen
% Informatik Klasse 7 | Stunde 3
% ═══════════════════════════════════════════════════════════════

\documentclass[11pt, a4paper]{article}

\usepackage[utf8]{inputenc}
\usepackage[T1]{fontenc}
\usepackage[ngerman]{babel}

\usepackage{helvet}
\renewcommand{\familydefault}{\sfdefault}

\usepackage[margin=1.5cm, top=1.8cm, bottom=1.5cm]{geometry}
\usepackage{enumitem}
\usepackage{tabularx}
\usepackage{booktabs}
\usepackage{array}
\usepackage{tikz}
\usepackage{fancyhdr}
\usepackage{xcolor}
\usepackage{tcolorbox}
\tcbuselibrary{skins}

% ═══════════════════════════════════════════════════════════════
% FARBEN
% ═══════════════════════════════════════════════════════════════

\definecolor{informatik}{HTML}{b35c00}
\definecolor{hellgrau}{HTML}{F5F5F5}
\definecolor{mittelgrau}{HTML}{888888}

\colorlet{fachfarbe}{informatik}

% ═══════════════════════════════════════════════════════════════
% KOPFZEILE
% ═══════════════════════════════════════════════════════════════

\pagestyle{fancy}
\fancyhf{}
\renewcommand{\headrulewidth}{0.4pt}
\renewcommand{\footrulewidth}{0pt}
\lhead{\small Informatik Klasse 7 | Tabellenkalkulation}
\rhead{\small Name: \underline{\hspace{3cm}}}

% ═══════════════════════════════════════════════════════════════
% BOXEN
% ═══════════════════════════════════════════════════════════════

\newtcolorbox{merkkasten}[1][]{
    colback=hellgrau,
    colframe=hellgrau,
    boxrule=0pt,
    arc=0pt,
    left=8pt, right=8pt, top=8pt, bottom=6pt,
    borderline north={2pt}{0pt}{fachfarbe},
    before skip=8pt, after skip=8pt,
    #1
}

\newtcolorbox{formelkasten}{
    colback=white,
    colframe=fachfarbe,
    boxrule=1pt,
    arc=0pt,
    left=10pt, right=10pt, top=8pt, bottom=8pt,
    before skip=6pt, after skip=6pt
}

\newtcolorbox{infokasten}{
    colback=white,
    colframe=gray!50,
    boxrule=1pt,
    arc=0pt,
    left=8pt, right=8pt, top=6pt, bottom=6pt,
    before skip=6pt, after skip=6pt
}

\setlength{\parindent}{0pt}
\setlength{\parskip}{4pt}

% ═══════════════════════════════════════════════════════════════
% DOKUMENT
% ═══════════════════════════════════════════════════════════════

\begin{document}

% Header
\begin{center}
    {\Large\bfseries\textcolor{fachfarbe}{Tabellenkalkulation Grundlagen}}\\[0.3em]
    {\small Lernblatt | Stunde 3}
\end{center}
\vspace{-0.3em}
\hrule
\vspace{0.8em}

% ═══════════════════════════════════════════════════════════════
% 1. GRUNDBEGRIFFE
% ═══════════════════════════════════════════════════════════════

\textbf{\textcolor{fachfarbe}{1. Grundbegriffe}}

\begin{merkkasten}
\textbf{Zelle:} Ein einzelnes Kästchen, in das man Daten (Text, Zahlen, Formeln) eingeben kann.
\end{merkkasten}

\begin{tabular}{@{}p{7cm}p{7cm}@{}}
\textbf{Spalte} & \textbf{Zeile} \\
$\bullet$ verläuft \textbf{senkrecht} (↓) & $\bullet$ verläuft \textbf{waagerecht} (→) \\
$\bullet$ wird mit \textbf{Buchstaben} bezeichnet & $\bullet$ wird mit \textbf{Nummern} bezeichnet \\
\hspace{0.5em} (A, B, C, D, ...) & \hspace{0.5em} (1, 2, 3, 4, ...)
\end{tabular}

\vspace{0.3em}
\textit{Merkhilfe:} \textbf{S}palte = \textbf{S}enkrecht

\vspace{0.5em}

% ═══════════════════════════════════════════════════════════════
% 2. ZELLADRESSEN
% ═══════════════════════════════════════════════════════════════

\textbf{\textcolor{fachfarbe}{2. Zelladressen}}

\begin{merkkasten}
Die \textbf{Zelladresse} besteht aus: \textbf{Spaltenbuchstabe + Zeilennummer}
\end{merkkasten}

\begin{infokasten}
\begin{center}
\begin{tikzpicture}[scale=0.7]
    % Header
    \fill[gray!30] (0,0) rectangle (1,0.6);
    \fill[gray!30] (1,0) rectangle (2.5,0.6);
    \fill[gray!30] (2.5,0) rectangle (4,0.6);
    \draw (0,0) rectangle (1,0.6);
    \draw (1,0) rectangle (2.5,0.6);
    \draw (2.5,0) rectangle (4,0.6);
    \node at (0.5,0.3) {\small};
    \node at (1.75,0.3) {\small\textbf{A}};
    \node at (3.25,0.3) {\small\textbf{B}};

    % Row 1
    \fill[gray!30] (0,-0.6) rectangle (1,0);
    \draw (0,-0.6) rectangle (1,0);
    \draw (1,-0.6) rectangle (2.5,0);
    \draw (2.5,-0.6) rectangle (4,0);
    \node at (0.5,-0.3) {\small\textbf{1}};
    \node at (1.75,-0.3) {\small A1};
    \node at (3.25,-0.3) {\small B1};

    % Row 2
    \fill[gray!30] (0,-1.2) rectangle (1,-0.6);
    \draw (0,-1.2) rectangle (1,-0.6);
    \draw (1,-1.2) rectangle (2.5,-0.6);
    \draw[fachfarbe, thick] (2.5,-1.2) rectangle (4,-0.6);
    \node at (0.5,-0.9) {\small\textbf{2}};
    \node at (1.75,-0.9) {\small A2};
    \node at (3.25,-0.9) {\small\textbf{B2}};

    % Pfeil
    \draw[->, thick, fachfarbe] (4.3,-0.9) -- (5.5,-0.9);
    \node[right] at (5.5,-0.9) {\small Spalte B, Zeile 2};
\end{tikzpicture}
\end{center}
\end{infokasten}

\vspace{0.3em}

% ═══════════════════════════════════════════════════════════════
% 3. FORMELN
% ═══════════════════════════════════════════════════════════════

\textbf{\textcolor{fachfarbe}{3. Formeln eingeben}}

\begin{merkkasten}
\textbf{Merke:} Jede Formel beginnt mit dem \textbf{Gleichheitszeichen} \texttt{=}
\end{merkkasten}

\textbf{Rechenzeichen in Formeln:}

\begin{center}
\begin{tabular}{c c c l}
\toprule
\textbf{Mathematik} & \textbf{Taste} & \textbf{Beispiel} & \textbf{Bedeutung} \\
\midrule
$+$ (plus) & \texttt{+} & \texttt{=A1+B1} & Addiert A1 und B1 \\
$-$ (minus) & \texttt{-} & \texttt{=A1-B1} & Subtrahiert B1 von A1 \\
$\times$ (mal) & \texttt{*} & \texttt{=A1*B1} & Multipliziert A1 mit B1 \\
$\div$ (geteilt) & \texttt{/} & \texttt{=A1/B1} & Dividiert A1 durch B1 \\
\bottomrule
\end{tabular}
\end{center}

\vspace{0.3em}

% ═══════════════════════════════════════════════════════════════
% 4. SUMME-FUNKTION
% ═══════════════════════════════════════════════════════════════

\textbf{\textcolor{fachfarbe}{4. Die SUMME-Funktion}}

\begin{formelkasten}
\centering
\large\texttt{=SUMME(Startadresse:Endadresse)}
\end{formelkasten}

\begin{itemize}[leftmargin=1.5em, itemsep=2pt]
    \item Addiert alle Werte in einem \textbf{Bereich}
    \item Der \textbf{Doppelpunkt} \texttt{:} bedeutet „von ... bis"
    \item Beispiel: \texttt{=SUMME(B2:B5)} addiert B2, B3, B4 \textbf{und} B5
\end{itemize}

\vspace{0.3em}

% ═══════════════════════════════════════════════════════════════
% 5. REFERENZPRINZIP
% ═══════════════════════════════════════════════════════════════

\textbf{\textcolor{fachfarbe}{5. Das Referenzprinzip}}

\begin{merkkasten}
\textbf{Referenzprinzip:} Ändert sich der Wert einer Zelle, werden alle Formeln, die diese Zelle verwenden, \textbf{automatisch neu berechnet}.
\end{merkkasten}

\begin{center}
\begin{tikzpicture}[scale=0.8]
    % Vorher
    \node[anchor=south] at (1.5,-0.6) {\small\textbf{Vorher:}};
    \draw (0,-1.2) rectangle (1.5,-0.6);
    \draw (1.5,-1.2) rectangle (3,-0.6);
    \node at (0.75,-0.9) {\small 10};
    \node at (2.25,-0.9) {\small 15};

    % Pfeil
    \draw[->, very thick, fachfarbe] (3.3,-0.9) -- (4.7,-0.9);
    \node[above] at (4,-0.9) {\tiny ändern};

    % Nachher
    \node[anchor=south] at (6.5,-0.6) {\small\textbf{Nachher:}};
    \draw[fachfarbe, thick] (5,-1.2) rectangle (6.5,-0.6);
    \draw (6.5,-1.2) rectangle (8,-0.6);
    \node at (5.75,-0.9) {\small\textbf{20}};
    \node at (7.25,-0.9) {\small\textbf{25}};

    % Beschriftung
    \node[below, font=\tiny] at (1.5,-1.3) {A1=10, B1=15};
    \node[below, font=\tiny] at (6.5,-1.3) {A1=20 → B1 wird 25};
\end{tikzpicture}
\end{center}

{\small In B1 steht \texttt{=A1+5}. Ändert man A1 auf 20, zeigt B1 automatisch 25.}

\vspace{0.5em}

% ═══════════════════════════════════════════════════════════════
% ZUSAMMENFASSUNG
% ═══════════════════════════════════════════════════════════════

\begin{tcolorbox}[colback=hellgrau, colframe=fachfarbe, boxrule=1pt, arc=0pt, title={\textbf{Zusammenfassung}}, fonttitle=\bfseries\color{white}, coltitle=white, colbacktitle=fachfarbe]
\begin{itemize}[leftmargin=1.5em, itemsep=2pt]
    \item \textbf{Zelle} = einzelnes Kästchen | \textbf{Adresse} = Spalte + Zeile (z.B. B2)
    \item \textbf{Spalte} = senkrecht (Buchstaben) | \textbf{Zeile} = waagerecht (Nummern)
    \item Formeln beginnen \textbf{immer} mit \texttt{=}
    \item Rechenzeichen: \texttt{+} \texttt{-} \texttt{*} \texttt{/}
    \item \texttt{=SUMME(B2:B5)} addiert einen Bereich (Doppelpunkt = „von...bis")
    \item \textbf{Referenzprinzip:} Automatische Neuberechnung bei Änderungen
\end{itemize}
\end{tcolorbox}

\end{document}
